% In this file you should put all LaTeX macros and settings to be used both by
% the pdf version and the web version.
% This should be most of your macros.

% Source-blueprint metadata consumed by hgraph and hidden in rendered prose.
\providecommand{\dcref}[1]{}
\providecommand{\uses}[1]{}
\providecommand{\proves}[1]{}
\providecommand{\lean}[1]{}
\providecommand{\leanok}{}
\providecommand{\mathlibok}{}
\providecommand{\notready}{}
\providecommand{\sketch}{}
\providecommand{\level}[1]{}

% The theorem-like environments defined below are those that appear by default
% in the dependency graph. See the README of leanblueprint if you need help to
% customize this.
% The configuration below use the theorem counter for all those environments
% (this is what the [theorem] arguments mean) and never resets it.
% If you want for instance to number them within chapters then you can add
% [chapter] at the end of the next line.
\newtheorem{theorem}{Theorem}
\newtheorem{proposition}[theorem]{Proposition}
\newtheorem{lemma}[theorem]{Lemma}
\newtheorem{corollary}[theorem]{Corollary}
\newtheorem{example}[theorem]{Example}

\theoremstyle{definition}
\newtheorem{definition}[theorem]{Definition}
\newtheorem{remark}[theorem]{Remark}

% Custom mathematical operators / macros used in the chapters.
% KaTeX (used by the dashboard) only recognises standard LaTeX/AMS macros, so
% non-standard commands must be defined here.
\DeclareMathOperator{\Spec}{Spec}
\DeclareMathOperator{\Pic}{Pic}
\DeclareMathOperator{\Ext}{Ext}
\DeclareMathOperator{\End}{End}
\DeclareMathOperator{\Sch}{Sch}
\DeclareMathOperator{\Hom}{Hom}
\DeclareMathOperator{\Sheaf}{Sheaf}
\DeclareMathOperator{\opens}{Op}
\DeclareMathOperator{\forget}{forget}
\DeclareMathOperator{\sheafCompose}{sheafCompose}
\DeclareMathOperator{\toAbSheaf}{toAbSheaf}
\DeclareMathOperator{\toModuleKSheaf}{toModuleKSheaf}
\DeclareMathOperator{\PicardFunctor}{PicardFunctor}
\DeclareMathOperator{\PicardFunctorAb}{PicardFunctorAb}
\newcommand{\Abelian}{\mathrm{Abelian}}
\newcommand{\AddCommGrpCat}{\mathrm{AddCommGrpCat}}
\newcommand{\CommRingCat}{\mathrm{CommRingCat}}
\newcommand{\RingCat}{\mathrm{RingCat}}
\newcommand{\ModuleCat}{\mathrm{ModuleCat}}
\newcommand{\Csp}{C_{\mathrm{top}}}
\newcommand{\struct}[1]{\mathcal{O}_{#1}}
\DeclareMathOperator{\finrank}{finrank}
\newcommand{\Module}{\mathrm{Module}}
\newcommand{\Linear}{\mathrm{Linear}}
\newcommand{\Algebra}{\mathrm{Algebra}}
\newcommand{\Over}{\mathrm{Over}}
\newcommand{\Opens}{\mathrm{Opens}}
\newcommand{\HasAffineCechAcyclicCover}{\mathtt{HasAffineCechAcyclicCover}}
\newcommand{\HasCechToHModuleIso}{\mathtt{HasCechToHModuleIso}}
\newcommand{\HasSheafify}{\mathtt{HasSheafify}}
\newcommand{\HasExt}{\mathtt{HasExt}}
\newcommand{\IsAffineHModuleVanishing}{\mathtt{IsAffineHModuleVanishing}}
\newcommand{\AddCommGroup}{\mathrm{AddCommGroup}}
\newcommand{\AlgCat}{\mathrm{AlgCat}}
\newcommand{\genus}{\mathrm{genus}}
\newcommand{\Jac}{\mathrm{Jac}}
\newcommand{\LineBundle}{\mathrm{LineBundle}}
\newcommand{\MonoidalCategory}{\mathrm{MonoidalCategory}}
\DeclareMathOperator{\HModule}{HModule}
\DeclareMathOperator{\HModulePrime}{HModule'}
\DeclareMathOperator{\HPrime}{H'}
\DeclareMathOperator{\cohomologyPresheaf}{cohomologyPresheaf}
\DeclareMathOperator{\extFunctor}{extFunctor}
\DeclareMathOperator{\presheafToSheaf}{presheafToSheaf}
\newcommand{\Functor}{\mathrm{Functor}}
\newcommand{\whiskeringRight}{\mathrm{whiskeringRight}}
\newcommand{\yoneda}{\mathrm{yoneda}}
\newcommand{\Category}{\mathrm{Category}}
\newcommand{\Scheme}{\mathrm{Scheme}}
\newcommand{\HasWeakSheafify}{\mathtt{HasWeakSheafify}}
\newcommand{\noncomputable}{\mathtt{noncomputable}}
\DeclareMathOperator{\free}{free}
\DeclareMathOperator{\toBiprod}{toBiprod}
\DeclareMathOperator{\fromBiprod}{fromBiprod}
\newcommand{\MayerVietorisSquare}{\mathrm{MayerVietorisSquare}}
\DeclareMathOperator{\shortComplex}{shortComplex}
\DeclareMathOperator{\isPushout}{IsPushout}
\DeclareMathOperator{\composeAndSheafify}{composeAndSheafify}
\DeclareMathOperator{\cokernelCofork}{cokernelCofork}
\DeclareMathOperator{\Exact}{Exact}
\DeclareMathOperator{\ShortExact}{ShortExact}
\newcommand{\PreservesMonomorphisms}{\mathrm{PreservesMonomorphisms}}
\newcommand{\Finsupp}{\mathrm{Finsupp}}
\newcommand{\CategoryTheory}{\mathrm{CategoryTheory}}
\newcommand{\toTopCat}{\mathrm{toTopCat}}
\newcommand{\IsAffineOpen}{\mathrm{IsAffineOpen}}
\newcommand{\AffineCoverMVSquare}{\mathrm{AffineCoverMVSquare}}
\newcommand{\toMayerVietorisSquare}{\mathrm{toMayerVietorisSquare}}
\newcommand{\mayerVietorisSquare}{\mathrm{mayerVietorisSquare}}
\newcommand{\grothendieckTopology}{\mathrm{grothendieckTopology}}
\newcommand{\cover}{\mathrm{cover}}
\newcommand{\Finite}{\mathrm{Finite}}
\newcommand{\toSquare}{\mathrm{toSquare}}
\newcommand{\HModulePrimeSeq}{\mathtt{HModule'\_sequence}}
\newcommand{\HModulePrimeSeqCurve}{\mathtt{HModule'\_sequence\_curve}}
\newcommand{\op}{\mathrm{op}}
\newcommand{\PUnit}{\mathrm{PUnit}}
\newcommand{\GrothendieckTopology}{\mathrm{GrothendieckTopology}}
\newcommand{\Type}{\mathrm{Type}}
\newcommand{\constantSheaf}{\mathrm{constantSheaf}}
\newcommand{\ULift}{\mathrm{ULift}}
\newcommand{\HModulePrimeTopSourceIso}{\mathtt{HModule'\_top\_sourceIso}}
\newcommand{\HModulePrimeTopLinearEquiv}{\mathtt{HModule'\_top\_linearEquiv}}
\DeclareMathOperator{\Sym}{Sym}
\DeclareMathOperator{\Div}{Div}
\DeclareMathOperator{\Alb}{Alb}
\newcommand{\Cech}{\check{C}}
\DeclareMathOperator{\rank}{rank}
% Component selectors used in Lean-flavoured prose (natural-transformation
% components, underlying morphism, projection arrows, underlying presheaf).
% \hom is provided rather than newly defined to avoid clashing with the
% AMS-LaTeX / KaTeX default \hom operator on the pdflatex side.
\newcommand{\app}{\mathrm{app}}
\providecommand{\hom}{\mathrm{hom}}
\newcommand{\pr}{\mathrm{pr}}
\newcommand{\presheaf}{\mathtt{presheaf}}

% plasTeX (leanblueprint web renderer) does not know the cleveref
% package; \cref{...} is otherwise rendered as `??<label>` in the web
% output. We fall back to \ref so the cross-references resolve. On
% pdflatex the cleveref package's real \cref is loaded later and takes
% precedence (this fallback is silently overridden by cleveref via
% \renewcommand-on-load).
\providecommand{\cref}[1]{\ref{#1}}
\providecommand{\Cref}[1]{\ref{#1}}
\providecommand{\crefrange}[2]{\ref{#1}--\ref{#2}}
\providecommand{\Crefrange}[2]{\ref{#1}--\ref{#2}}
% ── Macros referenced by chapters but previously undefined ──────────────
% (Found by a KaTeX render scan; without these both the dashboard and the
% plasTeX/leanblueprint build render the raw control sequences.)
\providecommand{\et}{\mathrm{\acute{e}t}}
\providecommand{\restriction}{\!\upharpoonright\!}
\providecommand{\llbracket}{[\![}
\providecommand{\rrbracket}{]\!]}
\providecommand{\zar}{\mathrm{Zar}}
\providecommand{\Set}{\mathrm{Set}}
\providecommand{\Sets}{\mathrm{Sets}}
\providecommand{\Ab}{\mathrm{Ab}}
\providecommand{\Sh}{\mathrm{Sh}}
\providecommand{\Psh}{\mathrm{PSh}}
\providecommand{\pp}{\mathfrak{p}}
\providecommand{\SheafHom}{\mathcal{H}\!om}
\providecommand{\GrpObj}{\mathrm{GrpObj}}
\providecommand{\F}{\mathcal{F}}
\providecommand{\OO}{\mathcal{O}}
\providecommand{\PP}{\mathbb{P}}
\providecommand{\Q}{\mathbb{Q}}
\providecommand{\Hilb}{\operatorname{Hilb}}
\providecommand{\Quot}{\operatorname{Quot}}
\providecommand{\Proj}{\operatorname{Proj}}
\providecommand{\Cl}{\operatorname{Cl}}
\providecommand{\ord}{\operatorname{ord}}
\providecommand{\val}{\operatorname{val}}
\providecommand{\id}{\operatorname{id}}
\providecommand{\Z}{\mathbb{Z}}
\providecommand{\fppf}{\mathrm{fppf}}
\providecommand{\red}{\mathrm{red}}
